\clearpage\thispagestyle{empty}
\pdfbookmark[0]{Quellen zu Pallasbeobachtungen (Oktober 1827)}{pallas1827-sources}
\begin{center}\textbf{Quellen zu Pallasbeobachtungen (Oktober 1827)}\end{center}
\noindent Diese KI-Übertragung ins heutige Deutsch folgt \textit{Pallas} in \href{https://gdz.sub.uni-goettingen.de/id/PPN236007467}{\textit{Werke}, Band~VI} (Göttingen, 1874), S.~460. Heinrich Christian Schumacher gibt die vier Beobachtungen wieder, die Gauß ihm mitgeteilt hat. Der Erstdruck erschien als \textit{Pallasbeobachtungen} im Oktober 1827 in den \textit{Astronomischen Nachrichten}, Band~VI, Nr.~124, \href{https://api.digitale-sammlungen.de/iiif/image/v2/bsb10049431_00050/full/1550,/0/default.jpg}{Sp.~67--68}. Der gebundene Band der Bayerischen Staatsbibliothek trägt die Jahreszahl 1828.
\par\smallskip\noindent Die Werkausgabe setzt die Einheiten in die erste Zeile und gibt Schumachers Nachnamen wieder, wo der Erstdruck mit \textit{S.} zeichnet. Die Tabelle und die Unterschrift folgen der Werkausgabe. KI-Übertragung und KI-Satz.
\clearpage
